% ECN 594: Midterm Review Notes
% This document contains review material and practice problems for the midterm exam
\documentclass[11pt]{article}
\usepackage{fullpage}
\usepackage[left=1.0in,top=1.0in,right=1.0in,bottom=1.0in]{geometry}
\usepackage{amsmath}
\usepackage{amssymb}
\usepackage{booktabs}
\usepackage{hyperref}

\usepackage[dvipsnames]{xcolor}
\definecolor{asumaroon}{rgb}{0.549,0.114,0.251}
\hypersetup{colorlinks,breaklinks,linkcolor=asumaroon,urlcolor=asumaroon,anchorcolor=asumaroon,citecolor=black}

\title{ECN 594: Midterm Review Notes}
\author{Nicholas Vreugdenhil}
\date{}

\begin{document}
\maketitle

\section{Demand Estimation}

\subsection{Why Demand Estimation?}
\begin{itemize}
	\item We need demand models to:
	\begin{itemize}
		\item Measure substitution patterns between products
		\item Compute price elasticities
		\item Evaluate policy (mergers, new products, price changes)
		\item Calculate consumer welfare
	\end{itemize}
	\item The dimensionality problem: $J$ products $\rightarrow$ $J^2$ elasticities
	\item Solution: characteristics-based models (Lancaster, BLP)
\end{itemize}

\subsection{The Logit Model: Key Equations}
\begin{itemize}
	\item Utility: $u_{ijt} = \delta_{jt} + \varepsilon_{ijt}$ where $\delta_{jt} = x_{jt}\beta + \alpha p_{jt} + \xi_{jt}$
	\item Choice probability (market share):
	\[
		s_{jt} = \frac{\exp(\delta_{jt})}{1 + \sum_{k=1}^J \exp(\delta_{kt})}
	\]
	\item Outside option: $s_{0t} = \frac{1}{1 + \sum_{k=1}^J \exp(\delta_{kt})}$
	\item \textbf{Berry inversion:}
	\[
		\ln(s_{jt}) - \ln(s_{0t}) = \delta_{jt}
	\]
\end{itemize}

\subsection{Logit Elasticities}
\begin{itemize}
	\item \textbf{Own-price elasticity:}
	\[
		\eta_{jjt} = \frac{\partial s_{jt}}{\partial p_{jt}} \cdot \frac{p_{jt}}{s_{jt}} = \alpha p_{jt} (1 - s_{jt})
	\]
	\item \textbf{Cross-price elasticity:}
	\[
		\eta_{jkt} = \frac{\partial s_{jt}}{\partial p_{kt}} \cdot \frac{p_{kt}}{s_{jt}} = -\alpha p_{kt} s_{kt}
	\]
	\item Note: $\alpha < 0$ so own-price elasticity is negative
	\item \textbf{IIA problem:} Cross-elasticity depends only on $k$, not on similarity to $j$
\end{itemize}

\subsection{The Identification Problem}
\begin{itemize}
	\item \textbf{Problem:} We observe equilibrium $(p, q)$ pairs
	\item Can't tell if demand shifted or supply shifted
	\item \textbf{Price endogeneity:} High unobserved quality $\xi_{jt} \rightarrow$ high price
	\begin{itemize}
		\item OLS sees: high price, still high demand
		\item Concludes: price doesn't matter much
		\item Result: $\hat{\alpha}$ biased toward zero
	\end{itemize}
	\item \textbf{Solution:} Instrumental variables
	\begin{itemize}
		\item Need: correlated with price, uncorrelated with $\xi$
		\item Examples: Hausman IVs, BLP IVs, cost shifters
	\end{itemize}
\end{itemize}

\subsection{Consumer Surplus: Log-Sum Formula}
\begin{itemize}
	\item Expected utility for consumer $i$:
	\[
		E[\max_j u_{ijt}] = \ln\left[\sum_{j=0}^J \exp(\delta_{jt})\right] + \text{constant}
	\]
	\item Consumer surplus (in dollars):
	\[
		CS = \frac{1}{|\alpha|} \ln\left[\sum_{j=0}^J \exp(\delta_{jt})\right]
	\]
	\item Divide by $|\alpha|$ to convert utils to dollars
	\item \textbf{Application:} Welfare change from adding/removing products
\end{itemize}

\subsection{The IIA Problem}
\begin{itemize}
	\item \textbf{IIA:} Ratio of choice probabilities doesn't depend on other options
	\[
		\frac{s_{jt}}{s_{kt}} = \frac{\exp(\delta_{jt})}{\exp(\delta_{kt})} = \exp(\delta_{jt} - \delta_{kt})
	\]
	\item \textbf{Red Bus / Blue Bus:} Adding identical bus option incorrectly increases welfare
	\item \textbf{Problem:} Logit doesn't know similar products are close substitutes
	\item \textbf{Solutions:}
	\begin{itemize}
		\item Demographic interactions (partial fix)
		\item Mixed logit / random coefficients (full fix, beyond scope)
	\end{itemize}
\end{itemize}

\subsection{Quick Summary: Logit Formulas}
\begin{center}
\begin{tabular}{ll}
	\toprule
	\textbf{Formula} & \textbf{Expression} \\
	\midrule
	Own-price elasticity & $\eta_{jjt} = \alpha p_{jt} (1 - s_{jt})$ \\
	Cross-price elasticity & $\eta_{jkt} = -\alpha p_{kt} s_{kt}$ \\
	Berry inversion & $\delta_{jt} = \ln(s_{jt}) - \ln(s_{0t})$ \\
	Consumer surplus & $CS = \frac{1}{|\alpha|} \ln\left[\sum_j \exp(\delta_{jt})\right]$ \\
	\bottomrule
\end{tabular}
\end{center}
Remember: $\alpha < 0$, so own-price elasticity is negative.

%%%%%%%%%%%%%%%%%%%%%%%%%%%%%%%%%%%%%%%%%%%%%%%%%%%%%%%%%%%%%
\section{Practice Problems: Demand Estimation}
%%%%%%%%%%%%%%%%%%%%%%%%%%%%%%%%%%%%%%%%%%%%%%%%%%%%%%%%%%%%%

\subsection{Practice: Elasticity Calculation}
\textbf{Question:} Suppose $\alpha = -0.5$, product A has price $p_A = 20$ and share $s_A = 0.1$.
\begin{itemize}
	\item (a) Calculate the own-price elasticity of product A.
	\item (b) If product B has $p_B = 25$ and $s_B = 0.15$, what is the cross-price elasticity of A with respect to B's price?
\end{itemize}

\textbf{Solution:}
\begin{itemize}
	\item \textbf{(a) Own-price elasticity:}
	\[
		\eta_{AA} = \alpha p_A (1 - s_A) = (-0.5)(20)(1 - 0.1) = -9
	\]
	Demand is elastic (a 1\% price increase reduces quantity by 9\%)
	\item \textbf{(b) Cross-price elasticity:}
	\[
		\eta_{AB} = -\alpha p_B s_B = -(-0.5)(25)(0.15) = 1.875
	\]
	A 1\% increase in B's price increases A's share by 1.875\%
\end{itemize}

\subsection{Practice: Identification (True/False/NEI)}
\begin{itemize}
	\item (a) OLS estimation of demand typically underestimates the price coefficient (in absolute value).
	\item (b) Using prices of the same product in other geographic markets as an IV is valid because prices in other markets don't affect local demand.
	\item (c) The logit model solves the dimensionality problem by assuming all products are equally substitutable.
\end{itemize}

\textbf{Solutions:}
\begin{itemize}
	\item \textbf{(a) TRUE.} Price endogeneity biases $\alpha$ toward zero. Since $\alpha < 0$, this means $|\hat{\alpha}_{OLS}| < |\alpha_{true}|$.
	\item \textbf{(b) TRUE (with caveat).} Hausman IVs work because common cost shocks affect prices in all markets (relevance), but other markets' prices don't directly affect local demand (exclusion). Caveat: requires no common demand shocks.
	\item \textbf{(c) FALSE.} Logit solves dimensionality by using product characteristics. But the IIA problem means substitution is proportional to share, not similarity.
\end{itemize}

\subsection{Practice: Berry Inversion}
\textbf{Question:} A market has 3 products plus outside option. Observed shares: $s_{0t} = 0.4$, $s_1 = 0.3$, $s_2 = 0.2$, $s_3 = 0.1$
\begin{itemize}
	\item (a) Compute the mean utilities $\delta_1$, $\delta_2$, $\delta_3$.
	\item (b) Rank the products by consumer preference.
\end{itemize}

\textbf{Solution:}
Using $\delta_{jt} = \ln(s_{jt}) - \ln(s_{0t})$:
\begin{align*}
	\delta_1 &= \ln(0.3) - \ln(0.4) = \ln(0.75) \approx -0.29 \\
	\delta_2 &= \ln(0.2) - \ln(0.4) = \ln(0.5) \approx -0.69 \\
	\delta_3 &= \ln(0.1) - \ln(0.4) = \ln(0.25) \approx -1.39
\end{align*}
\textbf{(b)} Ranking: Product 1 $>$ Product 2 $>$ Product 3. All products have negative mean utility relative to outside option (which is normalized to $\delta_0 = 0$).

\subsection{Practice: Consumer Surplus Change}
\textbf{Question:} Market has $\delta_0 = 0$, $\delta_1 = 1$, $\delta_2 = 0.5$. Price coefficient $\alpha = -0.5$. If product 1 is removed from the market, what is the consumer surplus loss per consumer?

\textbf{Solution:}
\begin{itemize}
	\item \textbf{Before removal:}
	\[
		CS^{\text{before}} = \frac{1}{0.5} \ln(e^0 + e^1 + e^{0.5}) = 2 \ln(1 + 2.72 + 1.65) = 2 \ln(5.37) \approx 3.36
	\]
	\item \textbf{After removal:}
	\[
		CS^{\text{after}} = \frac{1}{0.5} \ln(e^0 + e^{0.5}) = 2 \ln(1 + 1.65) = 2 \ln(2.65) \approx 1.95
	\]
	\item \textbf{CS Loss:} $3.36 - 1.95 = \$1.41$ per consumer
\end{itemize}

%%%%%%%%%%%%%%%%%%%%%%%%%%%%%%%%%%%%%%%%%%%%%%%%%%%%%%%%%%%%%
\section{Pricing and Price Discrimination}
%%%%%%%%%%%%%%%%%%%%%%%%%%%%%%%%%%%%%%%%%%%%%%%%%%%%%%%%%%%%%

\subsection{Monopoly Pricing: Key Equations}
\begin{itemize}
	\item Profit maximization: $\max_q \pi = p(q) \cdot q - c(q)$
	\item First-order condition: $MR = MC$
	\item \textbf{Lerner index:}
	\[
		L = \frac{p - MC}{p} = \frac{1}{|\varepsilon|}
	\]
	\item \textbf{Interpretation:} Markup is inversely related to elasticity
	\item Equivalently: $p = \frac{MC}{1 - 1/|\varepsilon|}$
\end{itemize}

\subsection{Types of Price Discrimination}
\begin{enumerate}
	\item \textbf{Perfect price discrimination}
	\begin{itemize}
		\item Charge each consumer their exact WTP
		\item Benchmark; rarely feasible
	\end{itemize}
	\item \textbf{Selection by indicators}
	\begin{itemize}
		\item Different prices for observable groups
		\item Example: student discounts, geographic pricing
	\end{itemize}
	\item \textbf{Self-selection}
	\begin{itemize}
		\item Design menu to induce type revelation
		\item Example: versioning, bundling
	\end{itemize}
\end{enumerate}

\subsection{Selection by Indicators: Inverse Elasticity Rule}
\begin{itemize}
	\item Two markets with different elasticities
	\item Apply Lerner in each market:
	\[
		\frac{p_1 - MC}{p_1} = \frac{1}{|\varepsilon_1|} \quad \text{and} \quad \frac{p_2 - MC}{p_2} = \frac{1}{|\varepsilon_2|}
	\]
	\item \textbf{Key insight:} Charge higher price in more inelastic market
	\item Rearranging: $p = \frac{MC}{1 - 1/|\varepsilon|}$
\end{itemize}

\subsection{Two-Part Tariffs}
\begin{itemize}
	\item Structure: $\text{Payment} = F + p \cdot q$
	\item \textbf{Optimal (homogeneous consumers):}
	\begin{itemize}
		\item Set $p = MC$ (maximize total surplus)
		\item Set $F = CS(p = MC)$ (extract all surplus)
	\end{itemize}
	\item \textbf{Result:} Firm captures entire surplus; efficient but inequitable
	\item \textbf{Heterogeneous consumers:} Tradeoff between extraction and participation
\end{itemize}

\subsection{Self-Selection: IC and IR Constraints}
\begin{itemize}
	\item \textbf{IC (Incentive Compatibility):} Each type prefers their option
	\[
		\text{IC}_H: \quad v_H^F - p_F \geq v_H^S - p_S
	\]
	\item \textbf{IR (Individual Rationality):} Each type willing to participate
	\[
		\text{IR}_L: \quad v_L^S - p_S \geq 0
	\]
	\item \textbf{Optimal menu:} IC binds for H, IR binds for L
	\item H gets ``information rent''; L gets zero surplus
\end{itemize}

%%%%%%%%%%%%%%%%%%%%%%%%%%%%%%%%%%%%%%%%%%%%%%%%%%%%%%%%%%%%%
\section{Practice Problems: Pricing}
%%%%%%%%%%%%%%%%%%%%%%%%%%%%%%%%%%%%%%%%%%%%%%%%%%%%%%%%%%%%%

\subsection{Practice: Lerner Index}
\textbf{Question:} A monopolist faces demand $p = 100 - 2q$ and has constant marginal cost $MC = 20$.
\begin{itemize}
	\item (a) Find the profit-maximizing price and quantity.
	\item (b) Calculate the Lerner index.
	\item (c) Verify using the elasticity formula.
\end{itemize}

\textbf{Solution:}
\begin{itemize}
	\item \textbf{(a)} $MR = 100 - 4q$. Set $MR = MC$: $100 - 4q = 20 \Rightarrow q = 20$. Price: $p = 100 - 2(20) = 60$
	\item \textbf{(b)} Lerner index: $L = \frac{60 - 20}{60} = \frac{40}{60} = \frac{2}{3}$
	\item \textbf{(c)} Elasticity: $\varepsilon = \frac{dq}{dp} \cdot \frac{p}{q} = -\frac{1}{2} \cdot \frac{60}{20} = -1.5$. Check: $L = \frac{1}{|\varepsilon|} = \frac{1}{1.5} = \frac{2}{3}$ $\checkmark$
\end{itemize}

\subsection{Practice: Selection by Indicators}
\textbf{Question:} A firm sells in two markets. Market A has $\varepsilon_A = -3$. Market B has $\varepsilon_B = -5$. Marginal cost is $MC = 10$. Find the optimal price in each market.

\textbf{Solution:}
Using $p = \frac{MC}{1 - 1/|\varepsilon|}$:
\begin{itemize}
	\item \textbf{Market A:} $p_A = \frac{10}{1 - 1/3} = \frac{10}{2/3} = 15$
	\item \textbf{Market B:} $p_B = \frac{10}{1 - 1/5} = \frac{10}{4/5} = 12.50$
\end{itemize}
Higher price in more inelastic market (A).

\subsection{Practice: Self-Selection}
\textbf{Question:} Two products (Premium, Basic), two consumer types.
\begin{center}
	\begin{tabular}{|c|c|c|}
		\hline
		& Premium & Basic \\
		\hline
		High type & 100 & 60 \\
		Low type & 50 & 40 \\
		\hline
	\end{tabular}
\end{center}
$MC = 20$ for both. Equal numbers of each type. A firm considers: $p_P = 100$, $p_B = 40$.
\begin{itemize}
	\item (a) What does each type buy?
	\item (b) Is this menu optimal? If not, what should change?
\end{itemize}

\textbf{Solution:}
\begin{itemize}
	\item \textbf{(a) Consumer choices:}
	\begin{itemize}
		\item High type: $CS_P = 100 - 100 = 0$, $CS_B = 60 - 40 = 20$. High type buys Basic! (IC violated)
		\item Low type: $CS_P = 50 - 100 = -50$, $CS_B = 40 - 40 = 0$. Low type buys Basic
	\end{itemize}
	\item \textbf{(b) Not optimal.} Need to lower $p_P$ so IC binds. IC binds when: $100 - p_P = 60 - 40 = 20 \Rightarrow p_P = 80$. Optimal menu: $p_P = 80$, $p_B = 40$
\end{itemize}

\subsection{Practice: Two-Part Tariff}
\textbf{Question:} A gym has demand $q = 20 - p$ from each consumer. Marginal cost is $MC = 4$.
\begin{itemize}
	\item (a) What is the optimal two-part tariff $(F, p)$?
	\item (b) What is the firm's profit per customer?
	\item (c) What would happen if firm charged $p = 10$ and $F = 50$?
\end{itemize}

\textbf{Solution:}
\begin{itemize}
	\item \textbf{(a)} Optimal: $p = MC = 4$. At $p = 4$: $q = 20 - 4 = 16$. Consumer surplus: $CS = \frac{1}{2}(20-4)(16) = 128$. Optimal fee: $F = 128$
	\item \textbf{(b)} Profit per customer: $F + (p - MC)q = 128 + 0 = 128$
	\item \textbf{(c)} At $p = 10$: $q = 10$, $CS = \frac{1}{2}(20-10)(10) = 50$. $F = 50$ is feasible, but profit = $50 + (10-4)(10) = 110 < 128$. Lesson: $p > MC$ leaves money on the table!
\end{itemize}

\subsection{Practice: Backing Out Marginal Cost}
\textbf{Question:} Firm observes price $p = 50$ and elasticity $\varepsilon = -2$.
\begin{itemize}
	\item (a) What is the firm's marginal cost?
	\item (b) If the elasticity were actually $-4$, what would $MC$ be?
	\item (c) What does this say about the importance of accurate elasticity estimates?
\end{itemize}

\textbf{Solution:}
Using Lerner: $MC = p\left(1 - \frac{1}{|\varepsilon|}\right)$
\begin{itemize}
	\item \textbf{(a)} $MC = 50(1 - 1/2) = 50(0.5) = 25$
	\item \textbf{(b)} $MC = 50(1 - 1/4) = 50(0.75) = 37.5$
	\item \textbf{(c)} Key insight: inferred MC is very sensitive to elasticity! Doubling $|\varepsilon|$ increased inferred $MC$ by 50\%. This matters for merger analysis: need accurate elasticities.
\end{itemize}

%%%%%%%%%%%%%%%%%%%%%%%%%%%%%%%%%%%%%%%%%%%%%%%%%%%%%%%%%%%%%
\section{Key Points Summary}
%%%%%%%%%%%%%%%%%%%%%%%%%%%%%%%%%%%%%%%%%%%%%%%%%%%%%%%%%%%%%

\begin{enumerate}
	\item \textbf{Logit:} $s_{jt} = \exp(\delta_{jt})/[1 + \Sigma\exp(\delta_{kt})]$
	\item \textbf{Berry inversion:} $\ln(s_{jt}) - \ln(s_{0t}) = \delta_{jt}$
	\item \textbf{Elasticities:} Own: $\alpha p_{jt}(1-s_{jt})$; Cross: $-\alpha p_{kt} s_{kt}$
	\item \textbf{IVs needed} because $E[p_{jt} \xi_{jt}] \neq 0$; bias toward zero
	\item \textbf{Lerner:} $L = (p-MC)/p = 1/|\varepsilon|$
	\item \textbf{Selection by indicators:} Higher price in inelastic market
	\item \textbf{Two-part tariff:} $p = MC$, $F = CS$
	\item \textbf{Self-selection:} IC binds for high type, IR for low type
\end{enumerate}

\end{document}
